\chapter{ComunicacionMotores.ino} \label{ComunicacionMotores.ino}
\blankpage

\begin{lstlisting}[language=Arduino, caption= ComunicacionMotores.ino, basicstyle=\scriptsize]
#include <mcp_can.h>
#include <SPI.h>

#ifdef ARDUINO_SAMD_VARIANT_COMPLIANCE
  #define SERIAL SerialUSB
#else
  #define SERIAL Serial
#endif

const int SPI_CS_PIN = 3; // MKR ?

MCP_CAN CAN(SPI_CS_PIN);

unsigned char Enter_Motor_Mode[8] = {0xFF, 0xFF, 0xFF, 0xFF, 0xFF, 0xFF, 0xFF, 0xFC};
unsigned char Exit_Motor_Mode[8] = {0xFF, 0xFF, 0xFF, 0xFF, 0xFF, 0xFF, 0xFF, 0xFD};

unsigned char buf;
unsigned char len = 0;
long unsigned int canID = 0;
unsigned char ID_type = 0;

void setup() {
  // put your setup code here, to run once:
  SERIAL.begin(115200);

}

void loop() {
  // put your main code here, to run repeatedly:

  while(CAN_OK != CAN.begin(0 ,CAN_1000KBPS, MCP_8MHZ)) {
    SERIAL.println("CAN BUS Shield init fail!");
    SERIAL.println("Init CAN BUS Shield again");
    delay(10);
  }

  SERIAL.println("CAN BUS Shield init ok!");

  CAN.sendMsgBuf(0x05, 0, 8, Enter_Motor_Mode);

  delay(1000);
  
  if(CAN_MSGAVAIL == CAN.checkReceive()){
    CAN.readMsgBuf(&canID, &ID_type, &len, &buf);
    SERIAL.println(len);
  
    SERIAL.println("----------------------");
    SERIAL.println("Get data from ID: 0x");
    SERIAL.println(canID, HEX);

    SERIAL.println(buf, HEX);
    SERIAL.println("\t");

    SERIAL.println();
  }

  else
    SERIAL.println("No llega nada");

  delay(1000);

}
\end{lstlisting}

\blankpage